% ==========================================================================
%  Chapter 66 — TrussElement: Embedded Rebar for Continuum Sub-Models
% ==========================================================================

\chapter{TrussElement: Embedded Rebar for Continuum Sub-Models}
\label{ch:truss-element}

% ==========================================================================
\section{Motivation}
\label{sec:truss-motivation}

The multiscale framework described in Chapter~\ref{ch:multiscale_seismic}
extracts prismatic sub-model volumes around beam--column joints that
exhibited yielding in the global structural analysis.  These sub-models
are discretised with hexahedral continuum elements and analysed with the
Ko--Bathe concrete model (Chapter~\ref{ch:ko-bathe-3d}).  However, a
realistic concrete sub-model must include steel reinforcement.

Rather than modelling rebar with continuum elements (which would require
extremely fine meshes), the standard approach in computational mechanics
is to embed one-dimensional \emph{truss} elements within the continuum
mesh.  When truss nodes coincide with hex element nodes, the two
element types share degrees of freedom through the DMPlex sieve
automatically---this is the \textbf{perfect-bond assumption}.

This chapter describes the implementation of a templated
\texttt{TrussElement<Dim>} that:
\begin{itemize}
    \item Models axial-only behaviour with a uniaxial constitutive law
          (e.g., Menegotto--Pinto steel),
    \item Integrates with the existing \texttt{FiniteElement} concept,
    \item Can be mixed with continuum elements in the same
          \texttt{Model} via \texttt{MultiElementPolicy} and \texttt{FEM\_Element},
    \item Supports both 2D and 3D spatial dimensions.
\end{itemize}


% ==========================================================================
\section{Formulation}
\label{sec:truss-formulation}

% --------------------------------------------------------------------------
\subsection{Kinematics}
\label{sec:truss-kinematics}

Consider a bar element connecting nodes $0$ and $1$ with position vectors
$\mathbf{x}_0$ and $\mathbf{x}_1$ in $\mathbb{R}^{\mathrm{Dim}}$.  Define:
\begin{equation}
    L = \|\mathbf{x}_1 - \mathbf{x}_0\|, \qquad
    \hat{\mathbf{e}} = \frac{\mathbf{x}_1 - \mathbf{x}_0}{L},
    \label{eq:truss-length-dir}
\end{equation}
where $L$ is the element length and $\hat{\mathbf{e}}$ the unit direction
vector.  For 3D, $\hat{\mathbf{e}} = (\ell, m, n)^T$ with $\ell = \cos\alpha_x$,
$m = \cos\alpha_y$, $n = \cos\alpha_z$ being the direction cosines.

The element has $2\,\texttt{Dim}$ translational DOFs arranged as:
\begin{equation}
    \mathbf{u}_e = (\mathbf{u}_0^T,\; \mathbf{u}_1^T)^T \in \mathbb{R}^{2\,\mathrm{Dim}}.
    \label{eq:truss-dof-vector}
\end{equation}

The axial strain is the projection of the displacement gradient onto the
element axis:
\begin{equation}
    \varepsilon = \hat{\mathbf{e}} \cdot
    \frac{\mathbf{u}_1 - \mathbf{u}_0}{L}
    = \mathbf{B}\,\mathbf{u}_e,
    \label{eq:truss-strain}
\end{equation}
where the $1 \times 2\,\mathrm{Dim}$ strain--displacement matrix is:
\begin{equation}
    \mathbf{B} = \frac{1}{L}
    \begin{pmatrix}
        -\hat{\mathbf{e}}^T & \hat{\mathbf{e}}^T
    \end{pmatrix}
    = \frac{1}{L}
    \begin{pmatrix}
        -\ell & -m & -n & \ell & m & n
    \end{pmatrix}.
    \label{eq:truss-B-matrix}
\end{equation}

Because the shape functions are linear, $\mathbf{B}$ is constant over the
element---there is no spatial variation.  A single Gauss point therefore
integrates all element integrals exactly.


% --------------------------------------------------------------------------
\subsection{Stiffness Matrix}
\label{sec:truss-stiffness}

The element stiffness follows from the virtual work principle.  With
cross-sectional area $A$ and elastic modulus $E$:
\begin{equation}
    \mathbf{K}_e = A \int_0^L \mathbf{B}^T E\, \mathbf{B}\;\mathrm{d}x
    = A \cdot L \cdot E \cdot \mathbf{B}^T \mathbf{B}.
    \label{eq:truss-K}
\end{equation}
The integration factor $w \cdot |J| = L$ comes from the single-point
Gauss rule on the reference interval $[-1,1]$ with weight $w=2$ and
Jacobian $|J| = L/2$.

Expanding $\mathbf{B}^T \mathbf{B}$ for 3D:
\begin{equation}
    \mathbf{K}_e = \frac{EA}{L}
    \begin{pmatrix}
        \phantom{-}\hat{\mathbf{e}}\hat{\mathbf{e}}^T
            & -\hat{\mathbf{e}}\hat{\mathbf{e}}^T \\[4pt]
        -\hat{\mathbf{e}}\hat{\mathbf{e}}^T
            & \phantom{-}\hat{\mathbf{e}}\hat{\mathbf{e}}^T
    \end{pmatrix},
    \label{eq:truss-K-expanded}
\end{equation}
where $\hat{\mathbf{e}}\hat{\mathbf{e}}^T$ is the $\mathrm{Dim}\times\mathrm{Dim}$
projection operator onto the element axis.  For an element aligned with the
$x$-axis ($\hat{\mathbf{e}} = (1,0,0)^T$), only the $K_{00}$, $K_{03}$,
$K_{30}$, $K_{33}$ entries are non-zero.


% --------------------------------------------------------------------------
\subsection{Internal Forces}
\label{sec:truss-internal-forces}

Given the current displacement $\mathbf{u}_e$, the internal force vector is:
\begin{equation}
    \mathbf{f}_e = A \cdot L \cdot \sigma(\varepsilon) \cdot \mathbf{B}^T,
    \label{eq:truss-f}
\end{equation}
where $\sigma(\varepsilon)$ is the uniaxial stress computed by the
constitutive model at the current strain $\varepsilon = \mathbf{B}\,\mathbf{u}_e$.


% --------------------------------------------------------------------------
\subsection{Tangent Stiffness}
\label{sec:truss-tangent}

The consistent tangent stiffness for nonlinear material behaviour:
\begin{equation}
    \mathbf{K}_t = A \cdot L \cdot E_t(\varepsilon) \cdot \mathbf{B}^T \mathbf{B},
    \label{eq:truss-Kt}
\end{equation}
where $E_t = \partial\sigma/\partial\varepsilon$ is the algorithmic tangent
modulus returned by the material model.


% --------------------------------------------------------------------------
\subsection{Mass Matrix}
\label{sec:truss-mass}

The consistent mass matrix for a 2-node bar with density $\rho$:
\begin{equation}
    \mathbf{M}_e = \frac{\rho A L}{6}
    \begin{pmatrix}
        2\mathbf{I} & \mathbf{I} \\
        \mathbf{I}  & 2\mathbf{I}
    \end{pmatrix},
    \label{eq:truss-M}
\end{equation}
where $\mathbf{I}$ is the $\mathrm{Dim} \times \mathrm{Dim}$ identity
matrix.  This gives lumped entries $\rho A L / 3$ on the diagonal
and coupling terms $\rho A L / 6$ on the off-diagonal blocks.


% ==========================================================================
\section{Implementation}
\label{sec:truss-implementation}

The element resides in a single header file,
\texttt{src/elements/TrussElement.hh} ($\sim$230~lines), and is templated
on the spatial dimension:

\begin{lstlisting}[style=cppstyle]
template <std::size_t Dim = 3>
class TrussElement {
    static_assert(Dim == 2 || Dim == 3);

    ElementGeometry<Dim>*      geometry_;
    Material<UniaxialMaterial> material_;  // 1 GP
    double                     area_;
    double                     L_;
    BMatrixT                   B_;        // precomputed (1 x 2*Dim)
    // ...
};
\end{lstlisting}

\noindent Key design decisions:

\begin{enumerate}
    \item \textbf{Precomputed B-matrix.}  Since $\mathbf{B}$ is constant
          for a 2-node element, it is computed once in the constructor from
          nodal coordinates and reused in all subsequent assembly calls.
          This avoids repeated geometry queries during nonlinear iterations.

    \item \textbf{Direct \texttt{Material<UniaxialMaterial>} storage.}  The
          element stores a type-erased uniaxial material handle
          directly, bypassing the \texttt{MaterialPoint} wrapper.  This is
          necessary because \texttt{MaterialPoint} binds a material to an
          \texttt{IntegrationPoint<dim>} where \texttt{dim} is the
          \emph{spatial} dimension---but the truss element lives in 3D
          space while its constitutive response is 1D.

    \item \textbf{DOF cache pattern.}  Following the same lazy-initialisation
          pattern as \texttt{ContinuumElement}, DOF indices are gathered from
          the \texttt{ElementGeometry} nodes on first use and cached for
          subsequent PETSc assembly calls.

    \item \textbf{Concept satisfaction.}  A \texttt{static\_assert} at
          file scope verifies that both \texttt{TrussElement<3>} and
          \texttt{TrussElement<2>} satisfy the \texttt{FiniteElement}
          concept (Section~\ref{sec:truss-concept}).
\end{enumerate}


% --------------------------------------------------------------------------
\subsection{FiniteElement Concept Compliance}
\label{sec:truss-concept}

The \texttt{FiniteElement} concept requires:
\begin{center}
\begin{tabular}{lcc}
    \toprule
    Method & ContinuumElement & TrussElement \\
    \midrule
    \texttt{num\_nodes()} & $n$ (8, 20, \ldots) & 2 \\
    \texttt{num\_integration\_points()} & $n_{gp}$ & 1 \\
    \texttt{sieve\_id()} & from geometry & from geometry \\
    \texttt{set\_num\_dof\_in\_nodes()} & $\mathrm{Dim}$ per node & $\mathrm{Dim}$ per node \\
    \texttt{inject\_K(Mat)} & $\int \mathbf{B}^T\mathbf{D}\mathbf{B}\,dV$
                             & $A L E\,\mathbf{B}^T\mathbf{B}$ \\
    \texttt{compute\_internal\_forces(Vec,Vec)} & $\int \mathbf{B}^T\bsigma\,dV$
                                                 & $A L \sigma\,\mathbf{B}^T$ \\
    \texttt{inject\_tangent\_stiffness(Vec,Mat)} & $\int \mathbf{B}^T\mathbf{D}_t\mathbf{B}\,dV$
                                                  & $A L E_t\,\mathbf{B}^T\mathbf{B}$ \\
    \texttt{commit\_material\_state(Vec)} & per GP & single GP \\
    \texttt{revert\_material\_state()} & per GP & single GP \\
    \bottomrule
\end{tabular}
\end{center}

The truss element assembles into the same global PETSc matrices via
\texttt{MatSetValuesLocal} and \texttt{VecSetValues}, interoperating
seamlessly with continuum elements in a mixed-topology model.


% --------------------------------------------------------------------------
\subsection{Integration with Mixed-Topology Models}
\label{sec:truss-mixed}

When truss elements coexist with hex elements in the same DMPlex mesh,
the model must use \texttt{MultiElementPolicy} (heterogeneous element
storage via type-erased \texttt{FEM\_Element} wrappers):

\begin{lstlisting}[style=cppstyle]
// Wrap both element types in FEM_Element
std::vector<FEM_Element> elements;
for (auto& hex : hex_elements)
    elements.emplace_back(std::move(hex));
for (auto& truss : truss_elements)
    elements.emplace_back(std::move(truss));

// Constructor 2: pre-built element vector
Model<MultiElementPolicy> model(domain, std::move(elements));
\end{lstlisting}

\noindent DOF coupling requires that truss nodes coincide with hex nodes
in the DMPlex sieve.  The \texttt{PrismaticDomainBuilder} generates the
hex mesh with node positions at regular grid points; rebar elements are
created by connecting pairs of these shared nodes along rebar paths.


% ==========================================================================
\section{Verification}
\label{sec:truss-verification}

A comprehensive unit test suite (\texttt{tests/test\_truss\_element.cpp},
$\sim$600~lines, 26~assertions) validates the implementation against
analytical solutions.

% --------------------------------------------------------------------------
\subsection{Test Material}
\label{sec:truss-test-material}

All tests use a Menegotto--Pinto steel model with very high yield stress
($f_y = 10^{12}$~MPa) to ensure purely elastic behaviour during
verification:
\begin{equation}
    E = 200\,000~\text{MPa}, \quad f_y = 10^{12}~\text{MPa}, \quad b = 0.01.
\end{equation}
Within the elastic range, $\sigma = E\varepsilon$ and $E_t = E$, giving
analytically predictable stiffness and force values.

% --------------------------------------------------------------------------
\subsection{Geometry Tests}
\label{sec:truss-geometry-tests}

\textbf{X-aligned bar} (nodes at $(0,0,0)$ and $(2,0,0)$):
\begin{itemize}
    \item \texttt{num\_nodes} = 2, \texttt{num\_integration\_points} = 1.
    \item $L = 2.0$.
\end{itemize}

\textbf{Diagonal bar} (nodes at $(0,0,0)$ and $(1,1,1)$):
\begin{itemize}
    \item $L = \sqrt{3}$.
\end{itemize}

% --------------------------------------------------------------------------
\subsection{Stiffness Tests}
\label{sec:truss-stiffness-tests}

\textbf{X-aligned bar} ($A=0.01$, $E=200\,000$, $L=2$):
\begin{equation}
    \frac{EA}{L} = \frac{200\,000 \times 0.01}{2} = 1000.
    \label{eq:truss-test-EA-L}
\end{equation}
Verified entries:
$K_{00} = +1000$,
$K_{03} = -1000$,
$K_{33} = +1000$,
$K_{11} = 0$ (transverse),
$K_{22} = 0$ (transverse).

\textbf{Diagonal bar} ($\hat{\mathbf{e}} = (1/\sqrt{3},\, 1/\sqrt{3},\, 0)^T$
for a bar from origin to $(1,1,0,0)$... actually $(1,1,1)$ with
$\cos^2\alpha = 1/3$):
\begin{equation}
    K_{00} = \frac{EA}{L}\cos^2\alpha = \frac{EA}{L}\cdot\frac{1}{3}, \quad
    K_{01} = \frac{EA}{L}\cos\alpha\cos\beta = \frac{EA}{L}\cdot\frac{1}{3}.
    \label{eq:truss-test-diagonal}
\end{equation}

% --------------------------------------------------------------------------
\subsection{Internal Force Test}
\label{sec:truss-force-test}

Imposing $\delta = 0.002$ axial elongation on the x-aligned bar:
\begin{equation}
    \varepsilon = \frac{\delta}{L} = 0.001, \quad
    \sigma = E\varepsilon = 200, \quad
    F = \sigma A = 2.0~\text{force units}.
\end{equation}
Verified: $f_0 = -F$, $f_3 = +F$, $f_1 = f_4 = 0$.

% --------------------------------------------------------------------------
\subsection{Additional Tests}

\begin{itemize}
    \item \textbf{Tangent consistency}: $\mathbf{K}_t(\mathbf{u}=\mathbf{0})
          = \mathbf{K}$ for elastic material (all $6 \times 6$ entries).
    \item \textbf{Commit/revert}: Material state after commit preserves
          strain; revert returns to last committed state.
    \item \textbf{FEM\_Element wrapping}: \texttt{TrussElement<3>} wraps
          successfully in \texttt{FEM\_Element} with preserved topology
          queries (\texttt{num\_nodes}, \texttt{num\_gp}, \texttt{sieve\_id}).
    \item \textbf{Mass matrix}: Diagonal entries $M_{00} = 2\rho AL/6$,
          off-diagonal $M_{03} = \rho AL/6$, cross-direction $M_{01} = 0$.
\end{itemize}

All 26 assertions pass.


% ==========================================================================
\section{Discussion}
\label{sec:truss-discussion}

The \texttt{TrussElement} completes the element infrastructure needed for
reinforced concrete sub-models.  With the concrete continuum modelled by
hexahedral elements using \texttt{KoBatheConcrete3D}
(Chapter~\ref{ch:ko-bathe-3d}) and steel reinforcement modelled by truss
elements using \texttt{MenegottoPintoSteel}, the sub-model solver can
now represent the composite behaviour of confined concrete with embedded
rebar.

The \textbf{perfect-bond assumption} (shared nodes) is appropriate for
well-confined regions typical of beam--column joints, where bond slip is
negligible under monotonic and low-cycle fatigue loading.

The next step is integrating the \texttt{TrussElement} into the
\texttt{SubModelSolver} (Chapter~\ref{ch:multiscale_seismic}), which
requires:
\begin{enumerate}
    \item A rebar mesh generator that creates line elements sharing nodes
          with the hex mesh inside the prismatic sub-model volume.
    \item Switching the model policy from \texttt{SingleElementPolicy} to
          \texttt{MultiElementPolicy}.
    \item Assigning distinct materials: \texttt{KoBatheConcrete3D} for hex
          elements and \texttt{MenegottoPintoSteel} for truss elements.
\end{enumerate}
